\section{推定と検定}
\label{b12_test}

\subsection{授業内容}
\subsubsection{概要}
これまでの授業では母集団分布と標本分布について学んできた。本コマでは標本から母集団の特性値(母数)を推定する方法について学ぶ。まず不偏推定量の概念を理解し，点推定と区間推定の違いについて学ぶ。標本統計量を用いて母数を推定する際の基本的な考え方と，推定値の精度や信頼性を評価する方法を理解する。さらに，推定と検定の関係についても触れ，統計的推論の基本的な枠組みを身につける。また，標準正規分布が統計的推論において果たす重要な役割についても理解を深める。

\subsubsection{コマ主題細目}
\begin{description}
	\item[不偏推定量] 母集団に正規分布を仮定し，標本平均の平均(期待値)が母平均に一致することを確認する。このことから，標本統計量が推定値として利用できることを確認する。ただし，分散の場合は一致しない(不偏性がない)ことをみる。分散については不偏分散を考えることでこの不一致を解決することができる。標本統計量が母数の良い推定量であるための条件(不偏性，一致性，効率性)についても理解する。また，推定量と推定値の違いについても理解する。
	      \begin{flushright}
		      $\to$ 不偏推定量については\textcite{Yamada2004}のP.98--103に詳しい。標本統計量のこれらの特徴については，\textcite{kosugi2023}のP.119--141も参照すると良い。
	      \end{flushright}

	\item[点推定と標準誤差] 点推定とは，標本統計量を用いて母数の値を単一の値として推定する方法である。例えば，標本平均は母平均の点推定値となる。しかし，この推定値にはサンプリング誤差が含まれるため，その精度を評価する必要がある。標準誤差は推定値の精度を表す重要な指標であり，標本平均の標準誤差は母標準偏差をサンプルサイズの平方根で割ったものとして定義される。標準誤差の意味と解釈，およびサンプルサイズとの関係について理解する。また，中心極限定理により，サンプルサイズが大きくなると標本平均の分布が正規分布に近づくことも確認する。
	      \begin{flushright}
		      $\to$ \textcite{Kawabata201408}のP.107--110，\textcite{Yamada2004}のP.116--119.
	      \end{flushright}

	\item[区間推定と信頼区間] 標本平均は母平均の不偏推定量ではあるが，ある値をそのまま母平均であると考えるのは(標準誤差の考え方があるとは言え)やや断定的である。これに対して一定の幅を使って予測する区間推定という方法がある。推定に確率の言葉が入ってくるので不慣れな表現が出てくるが，注意深く解釈する必要がある。95\%信頼区間の意味と解釈，標準誤差を用いた信頼区間の計算方法について理解する。また，信頼区間の幅とサンプルサイズの関係についても学ぶ。信頼区間の正しい解釈は，「この方法で区間を構成すると，長期的には95\%の確率で真の母数を含む」というものであり，特定の区間が母数を含む確率が95\%という意味ではないことに注意する。
	      \begin{flushright}
		      $\to$ \textcite{Kawabata201408}のP.111--116，\textcite{Yamada2004}のP.120--127.
	      \end{flushright}


	\item[推定と検定の関係] 推定は母数の値を推し量る作業であるのに対し，検定はある仮説の下での統計量の期待値と実際の統計量を比較することで，仮説の妥当性を判断する作業である。両者は密接に関連しており，推定は「どれくらいか」という問いに，検定は「差があるか否か」という問いに答えるものである。推定と検定の関係を理解することで，次回学ぶ帰無仮説検定の基礎を形成する。特に，信頼区間と検定の関係について理解を深める。
	      \begin{flushright}
		      $\to$ \textcite{Kawabata201408}のP.81--85，\textcite{Yamada2004}のP.104--110.
	      \end{flushright}

\end{description}
\subsubsection{キーワード}
\begin{itemize}
	\item 不偏推定量
	\item 点推定と標準誤差
	\item 区間推定と信頼区間
	\item 標準正規分布
	\item 推定と検定の関係
\end{itemize}
\subsection{授業情報}
\paragraph{コマの展開方法}
講義
\paragraph{標準シラバスにおける位置づけ}
科目番号5；心理学統計法；1.心理学で用いられる統計手法；F 推測統計:その考え方
\subsubsection{予習・復習課題}
\paragraph{予習}
前回までに学んだ標本分布，標準誤差，標準正規分布についての理解を確認しておくこと。特に，標本統計量の標準化の方法と標準正規分布表の読み方について復習しておくと理解が深まる。
\paragraph{復習}
授業で学んだ内容を定着させるために，以下の点について理解を深めておくこと：

\begin{enumerate}
\item 不偏推定量の概念と，標本分散と不偏分散の違い
\item 点推定と標準誤差の関係，およびサンプルサイズが推定精度に与える影響
\item 信頼区間の正しい解釈と，その構成方法
\item 標準正規分布を用いた統計的推論の基本的な考え方
\item 推定と検定の関係，および信頼区間と検定の関連性
\end{enumerate}

統計的推定は，実証的な心理学研究の基盤をなす重要な概念である。特に，点推定だけでなく区間推定の考え方を理解することで，データの持つ不確実性を適切に評価できるようになる。統計的推定についての理解を深めるには，\textcite{Kawabata201408}のP.104--116や\textcite{Yamada2004}のP.98--127などを参考にすると良い。